\section*{Chapitre 19 - Probabilités}
 \addcontentsline{toc}{section}{Chapitre 19 - Probabilités}
 
\subsection*{I. Vocabulaire}
\addcontentsline{toc}{subsection}{I. Vocabulaire}

\begin{Définitions*}[c]
\leavevmode\vspace{-\baselineskip}
\begin{itemize}[label=\textcolor{Couleur6}{$\bullet$}]
\item Une expérience est \textbf{aléatoire} lorsqu'on ne peut pas en prévoir le résultat.
\item Une \textbf{issue} est un résultat possible d'une expérience aléatoire.
\item Un \textbf{événement} d'une expérience aléatoire est un ensemble d'issues.
\end{itemize}
\end{Définitions*}

\vspace{0.3cm}

\begin{Exemple}  Sur un jeu de 13 cartes indiscernables au toucher, on écrit sur chaque carte une lettre du mot \\
\og mathématiques \fg{} : M A T H E M A T I Q U E S. \\
On retourne toutes les cartes et demande à une personne d'en piocher une au hasard.
\begin{itemize}[label=\textcolor{Couleur8}{$\bullet$}]
\item Cette expérience est aléatoire car on ne sait pas la lettre qui sera piochée (puisque les cartes sont retournées).
\item Les issues de l'expérience sont : M, A, T, H, E, I, Q, U, S
\item Un événement possible est : \og Tirer une voyelle \fg{}. Il y a plusieurs issues qui réalisent cet événement : \og A ; E ; I ; U \fg{}.
\end{itemize}
\end{Exemple}


\subsection*{II. Probabilité d'un événement}
\addcontentsline{toc}{subsection}{II. Probabilité d'un événement}

\vspace{0.3cm}

\begin{Définition}[c]
Dans une expérience aléatoire, la \textbf{probabilité} d'un évènement donne une mesure de la \og chance \fg{} que cet évènement se produise. \\La probabilité d'un évènement est un nombre compris entre 0 et 1.

\end{Définition}


\begin{center}
\begin{tikzpicture}[scale=0.95]
	\coordinate (A) at (0,0);
	\coordinate (B) at (4,0);
	\coordinate (C) at (8,0);
 	\coordinate (D) at (12,0);
	\coordinate (E) at (16,0); 	   


    \draw[line width=4pt, Couleur8] (A) -- (B);  
    \draw[line width=4pt, Couleur7] (B) -- (C);  
    \draw[line width=4pt, Couleur11] (C) -- (D);  
    \draw[line width=4pt, Couleur9] (D) -- (E);    
    \draw[line width=0.5pt,gray] (A) -- (E); 
    \draw[line width=1.5pt, Couleur12] (0,-0.2) -- (0,0.2);  
    \draw[line width=1.5pt, Couleur6] (8,-0.2) -- (8,0.2);  
    \draw[line width=1.5pt, Couleur4] (16,-0.2) -- (16,0.2);  


    
    
    % Labels des points
    \node[Couleur12,below] at (0,-0.25) {\textbf{0}};
    \node[Couleur6,below] at (8,-0.25) {\textbf{0,5}};
    \node[Couleur4,below] at (16,-0.25) {\textbf{1}};
    

    \node[Couleur12,above] at (0,0.3) {\textbf{Impossible}};
    \node[Couleur8,above] at (2,-1) {\textbf{Improbable}};
    \node[Couleur7,above] at (6,-1) {\textbf{Peu porbable}};
    \node[Couleur6,above] at (8,0.3) {\textbf{Une chance sur deux}};
    \node[Couleur11,above] at (10,-1) {\textbf{Probable}};
    \node[Couleur9,above] at (14,-1) {\textbf{Très probable}};
    \node[Couleur4,above] at (16,0.3) {\textbf{Certain}};


\end{tikzpicture}
\end{center}



\vspace{0.3cm}
\begin{Propriété}
Dans une situation d'équiprobabilité, la probabilité d'un évènement est la proportion du nombre d'issues favorables par rapport au nombre total d'issues possibles. \\
La probabilité est donc égale à la fraction : 
\begin{center}
$\dfrac{\text{Nombre d'issues favorables}}{\text{Nombre total d'issues possibles}}$
\end{center}


\end{Propriété}

\vspace{0.3cm}

\begin{Exemple}  On lance un dé équilibré à 6 faces. \\
Quelle est la probabilité d'obtenir un nombre supérieur ou égal à 5 ? \\
Cet événement possède 2 issues possibles sur 6 issues (\og 5 \fg{} et \og 6 \fg{}). \\
On a donc 2 chances sur 6, c'est-à-dire une probabilité égale à $\dfrac{2}{6}=\dfrac{1}{3}$.
\end{Exemple}

